% ============================================================
%  ch02c：§2.3 受限几何朗兰兹（原文 p.12）
% ============================================================
\subsection{受限几何朗兰兹}

\pnum{}令 $\Shv_{\Nilp}(\Bun_G) \subseteq \Shv(\Bun_G)$ 表示具有\emph{幂零锥中的奇异支撑}的层子范畴。这里 \et{} 层的奇异支撑由 Beilinson 在 [5] 中定义。

在 Betti 设定中，Ben-Zvi 与 Nadler [7] 说具有幂零奇异支撑的层是应当研究的正确对象（在 [2] 定理 18.1.6 中可以找到一个证明这一想法的精确定理）。我们把这一原则模仿到 \et{} 设定中，猜想：

\begin{conj}[受限几何朗兰兹猜想]
存在一个等价
\[
  \Shv_{\Nilp}(\Bun_G) \simeq \IndCoh_{\Nilp^{\mathrm{spec}}}(\LS^{\restr}_{\bG}).
\]
\end{conj}

这里右端如在 de Rham 情形中那样定义。

\begin{remark}
人们可以在 [23] 猜想 6.3.2 中找到这个猜想的一个简化版本。
\end{remark}

\begin{remark}
当然，猜想 2.3.1 受制于许多相容性。与 Whittaker 系数的相容性、（与 Eisenstein 级数的相容性的一个温和形式）\footnote{关于 Whittaker 相容性与 Eisenstein 相容性的表述，参见 [13]。} 以及 [10] 定理 8.3.0.1 的一个版本，\emph{唯一地}确定了猜想 2.3.1 中的比较函子；在 de Rham 与 Betti 设定中，这一想法是 [16] 的主题。
\end{remark}

\pnumt{证据。}{} 当几何域 $k$ 特征为零时，我们在 [2] 中证明受限 GLC 可由 de Rham 几何朗兰兹猜想推出。

一般地，人们可以直接对 $G = \mathbb{G}_m$ 验证该猜想；并且，用与 [21] 类似的想法，可以把该猜想归约到 $X = \mathbb{P}^1$ 的导出 Satake。

